% !TeX root = bash.tex
\section{I. File Tree}
\begin{frame}[fragile]
\frametitle{Everything is a "file"}
File basically means something that could be read or written. These are all treated as files:
\begin{itemize}
    \item \tt{a}
    \item \tt{.a} (Hidden)
    \item \tt{a.jpg}
    \item \tt{/dev/sda} (hardware)
    \item \tt{/dev/null} (trash bucket)
\end{itemize}

\begin{block}{Note}
    The dot and suffix are part of the filename.
    \newline
    Filenames are case \textbf{sensitive} in linux but case \textbf{insensitive} in macOS.
    \newline
    \textbf{Avoid spaces and special characters} (except \verb|._-|).
    If you have to, surround filename in quotes:
    \tt{`Lab Report (3).docx'}
\end{block}
\end{frame}

\begin{frame}
\frametitle{Directories}
Each of these is a \textbf{directory} (``dir'' for short):
\begin{itemize}
    \item \tt{01-files/}
    \item \tt{01-files/c/}
    \item \tt{01-files/.c/} (Hidden dir)
\end{itemize}

\begin{block}{Convention}
    For clarity, we add a slash (\tt{/}) to the end of a directory in the
    slides. However, in reality it often makes no difference.
\end{block}
\end{frame}

\begin{frame}
\frametitle{File tree}
Think of any directory as a tree.

\begin{figure}[h]
    \centering
    \begin{tikzpicture}
        \tikzstyle{every node}=[fill=blue!20,rounded corners]
        \tikzstyle{edge from parent}=[blue!50,thick,draw]
        \node {01-files/}[edge from parent fork down]
            child {node {a.txt}}
            child {
                node {c/}
                    child {node {subdir/}}
            }
            child {node {.c/}};

    \end{tikzpicture}
\end{figure}
\end{frame}

\begin{frame}
\frametitle{Paths}
    dir/subdir1/subdir2/file = \textbf{path}.
\footnote{At least in the scope of this workshop.}
\newline \newline

\begin{itemize}
    \item we use path to specify a file \& navigate the file system
    \item working directory is the path that you are currently in the file system
\end{itemize}
\end{frame}

\begin{frame}
\frametitle{Absolute \& relative paths}
\begin{itemize}
    \item \textbf{/} is the \textbf{root} directory of the file system. Paths beginning with \tt{/} are absolute: \tt{/usr/bin/cat}
    \item Otherwise it is relative: \tt{01-files/} (relative to working directory)
\end{itemize}


\begin{example}
    Your working directory: \tt{/home/you/} \newline
    Absolute path: \tt{/home/you/bash-workshop/01-files/} \newline
    Relative path(the same): \tt{bash-workshop/01-files/} \newline
\end{example}
\end{frame}

\begin{frame}
\frametitle{\tt{.} and \tt{..}}
Inside every dir\footnote{Except \tt{/}} there are two special dirs:
\begin{itemize}
    \item \tt{./} — current dir
    \item \tt{../} — parent dir
\end{itemize}

You can use them in relative paths.
\begin{example}
    Your location: \tt{/home/you/} \newline
    The followings are equal:
    \begin{itemize}
        \item \tt{./project/01-files}
        \item \tt{project/../project/01-files}
    \end{itemize}
\end{example}
\end{frame}

\begin{frame}[fragile]
\frametitle{Some handy tools/commands}
\begin{columns}
\begin{column}{0.5\textwidth}
\textbf{External tools come with OS}:
\begin{itemize}
    \item \tt{cp} — copy files
    \item \tt{mv} — move/rename files
    \item \tt{rm} — remove files
    \item \tt{ls} — list directory contents
    \item \tt{cat} — print file contents
\end{itemize}
\end{column}
\begin{column}{0.5\textwidth}
\textbf{Bash built-in commands}:
\begin{itemize}
    \item \tt{pwd} — print working directory
    \item \tt{cd} — change directory
    \item \tt{echo} — print text
    \item \tt{type} — show command type
\end{itemize}
\end{column}
\end{columns}
\vspace{0.5em}
\begin{block}{Tip}
    Use \tt{type <cmd>} to check if a command is built-in or external.
\end{block}
\end{frame}

\begin{frame}[fragile]
\frametitle{\tt{cat}: Printing a file}
Open a bash terminal inside \tt{01-files/}, then:
\begin{lstlisting}[language=bash]
$ cat a
$ cat .a
$ cat a.txt
\end{lstlisting}
\begin{block}{Explanation}
    \tt{cat} is short for ``concatenate'' (to join together) but it's
    mostly used to print files.
\end{block}
\end{frame}

\begin{frame}[fragile]
\frametitle{\tt{cp, mv, rm}: Relocating a file}
Try this inside \tt{01-files/}:
\begin{lstlisting}[language=bash]
$ ls
$ cp a b
$ ls
$ mv a.txt b.txt
$ ls
$ rm b
$ ls
\end{lstlisting}
\begin{block}{Explanation}
    \begin{itemize}
        \item \textbf{ls} lists files
        \item \textbf{cp} copies \tt{a} into a file called \tt{b}
        \item \textbf{mv} moves \tt{a.txt} into a file called \tt{b.txt}
        \item \textbf{rm} removes \tt{b} (\textbf{Caution! Irreversible})
    \end{itemize}
\end{block}
\end{frame}

\begin{frame}[fragile]
\frametitle{\tt{cp, mv}: Overwriting and renaming}
When the destination does not exist, \tt{cp} and \tt{mv} simply create
that file. \textbf{Otherwise, it is destroyed and overwritten.}
\newline \newline
Try this inside \tt{01-files/}:
\begin{lstlisting}[language=bash]
$ cp a b      # creates b
$ cat a
$ cp a a.txt  # overwrites a.txt
$ cat a.txt
\end{lstlisting}

Renaming a file in bash works like so:
\begin{lstlisting}[language=bash]
$ mv b newb
\end{lstlisting}
\end{frame}


\begin{frame}[fragile]
\frametitle{\tt{cd, pwd}: Changing directory}
Try this inside \tt{01-files/}:
\begin{lstlisting}[language=bash]
$ cd c/
$ pwd
$ cd ../
$ pwd
\end{lstlisting}
\begin{block}{Explanation}
    \begin{itemize}
        \item \tt{cd}: ``change directory''
        \item \tt{pwd}: ``print working directory''
        \item \tt{../} means ``parent directory''
    \end{itemize}
\end{block}
\end{frame}

\begin{frame}[fragile]
\frametitle{\tt{ls}: Listing directories}
Try this inside \tt{01-files/}:
\begin{lstlisting}[language=bash]
$ ls
$ ls -a
$ ls -l
$ ls -la
$ ls c/
\end{lstlisting}
\begin{block}{Explanation}
    \begin{itemize}
        \item \tt{ls}: ``list''
        \item \tt{-la} = \tt{-l} + \tt{-a}. Sth like this are called a `flag'
        \item \tt{-a} is short for \tt{--all}
        \item \tt{-l} enables long listing format
    \end{itemize}
\end{block}
\end{frame}

\begin{frame}[fragile]
\frametitle{\tt{mkdir, rm}: Creating and deleting directories}
Try this inside \tt{01-files/}:
\begin{lstlisting}[language=bash]
$ mkdir dir/
$ rm -r dir/
$ mkdir -p dir/subdir/
$ ls dir/
\end{lstlisting}
\begin{block}{Explanation}
    \begin{itemize}
        \item \tt{mkdir}: ``make directory''
        \item \tt{-r} is short for \tt{--recursive}
        \item \tt{-p} is short for \tt{--parents}. It makes sure the parent directory exists.  
    \end{itemize}
\end{block}
\end{frame}

\begin{frame}[fragile]
\frametitle{\tt{cp, mv}: Into and out of directories}
Try this inside \tt{01-files/}:
\begin{lstlisting}[language=bash]
$ cp a dir/
$ mv b.txt dir/
$ mv dir/a b.txt
\end{lstlisting}
\begin{block}{Explanation}
    \begin{itemize}
        \item Copy \tt{a} into \tt{dir/}
        \item Move \tt{b.txt} into \tt{dir/}
        \item Move \tt{dir/a} back into \tt{b.txt}
    \end{itemize}
\end{block}
\end{frame}


\begin{frame}
\frametitle{Challenge}
Inside \tt{01-files/}:
\begin{itemize}
    \item Enter \tt{challenge/}
    \item Create \tt{backup/}
    \item Copy \tt{a.txt} into \tt{dir/}
    \item Move \tt{dir/} into \tt{backup/}
    \item Verify using \tt{ls}
    \item Delete \tt{backup/}
    \item Go to the next section's directory
\end{itemize}
\end{frame}

\begin{frame}[fragile]
\frametitle{Solution}
\begin{lstlisting}[language=bash]
$ cd challenge/
$ mkdir backup/
$ cp a.txt dir/
$ mv dir/ backup/
$ ls
# Output: backup/  a.txt
$ ls backup/dir/
# Output: a.txt
$ rm -r backup/
$ cd ../../02-cli/
\end{lstlisting}
\end{frame}
